\scanpage{279}
\mspage{20}

Il est bon de regarder \deleted{\unsure{les}} le \uline{type} de chemins comme une donnée \deleted{\unsure{d'un}} non catégorique, mais combinatoire, une simple suite $\tau$ de $n$ signes \inserted{$\tau_i$, $(1\leq i\leq n)$} $+$ ou $-$ ($n={}$ longueur des chemins envisagés). \deleted{\unsure{Et}} Cette donnée combinatoire définit un intervalle orienté \unsure{standard|standart}, \deleted{qui \unsure{pa}} $I_\tau$.

Quand on a un intervalle orienté quelconque $(I,\omega)$, il définit un type $\tau$, et on l'identifie canoniquement au type $I_\tau$. Il \deleted{\unsure{vaut}} vaut mieux écrire $\underline{\mathrm{Ch}}_\tau(X)$, \unsure{plus tôt|plutôt} que $\underline{\mathrm{Ch}}_{\unsure{(}I,\omega)}(X)$, pour bien indiquer que pour $X$ donné, c'est là une catégorie totalement ``épinglée'', \unsure{sans|des} automorphismes liés à la nature particulière de $\tau$ (\unsure{ô|à} les \unsure{autom.} de $\tau$ !).

On peut préciser ce point ainsi. Soit $\tau'$ un autre type d'intervalles, et proposons-nous de déterminer des \deleted{\unsure{fonct}} ``lois'' qui permettent, à tout \deleted{\unsure{che}} chemin de type $\tau$ dans une cat. \inserted{arbitraire} $X$, d'associer un chemin de type $\tau'$ dans cette \unsure{$\overline{\mathrm m}$|$=$} catégorie, et ceci de façon \inserted{``fonctorielle'', i.e.} compatible avec les \deleted{foncteurs} applications horizontales
\diagram{\begin{tikzcd}[ampersand replacement=\&,row sep=small]\mathrm{Ch}_\tau(X) \arrow[r,"{\mathrm{Ch}_\tau(f)}"] \arrow[d,"{L_X}"'] \& \mathrm{Ch}_\tau(X') \arrow[d,"{L_{X'}}"]\\
\mathrm{Ch}_{\tau'}(X) \arrow[r,"{\mathrm{Ch}_{\tau'}(f)}"] \&\mathrm{Ch}_{\tau'}(X')
\end{tikzcd}}
associées \deleted{aux foncteurs \unsure{q|p}} \inserted{à un} foncteur $f\colon X\mathrel{\unsure{\text{---}|\longrightarrow}}X'$.
% archive sheet number: 278
